% =====================================================================
%  ANNEXE A — Évaluation multi-LLM du mapping MITRE ATT&CK
%             (workflow n8n sans MCP) et calcul des métriques
% ---------------------------------------------------------------------
%  Autonome et compilable avec pdflatex. Pour l'intégrer à un mémoire,
%  copier le contenu entre \begin{document} et \end{document} après
%  \appendix, et reporter les paquets « REQUIS » dans le préambule.
% =====================================================================
\documentclass[11pt,a4paper]{article}

% ---------- Paquets REQUIS ----------
\usepackage[utf8]{inputenc}
\usepackage[T1]{fontenc}
\usepackage[main=french,provide=*]{babel}
\IfFileExists{lmodern.sty}{\usepackage{lmodern}}{}
\usepackage[margin=1.9cm]{geometry}
\usepackage{xcolor}
\usepackage{amsmath}
\usepackage{array}
\usepackage{booktabs}
\usepackage{caption}
\usepackage{listings}
\usepackage{tikz}
\usetikzlibrary{positioning,arrows.meta,calc,fit,backgrounds,shapes.geometric}

% ---------- Palette ----------
\definecolor{cIn}{HTML}{2C3E50}\definecolor{cProc}{HTML}{2980B9}
\definecolor{cDec}{HTML}{E67E22}\definecolor{cLLM}{HTML}{8E44AD}
\definecolor{cGT}{HTML}{27AE60}\definecolor{cErr}{HTML}{C0392B}
\definecolor{cGy}{HTML}{7F8C8D}
\definecolor{cBg}{HTML}{F4F6F7}\definecolor{cCode}{HTML}{ECF0F1}
\definecolor{cComment}{HTML}{7F8C8D}\definecolor{cStr}{HTML}{C0392B}

% ---------- Styles TikZ ----------
\tikzset{
  bul/.style={rounded corners=4pt, draw, thin, align=center, font=\tiny,
              inner sep=2pt, minimum height=0.5cm, text width=2.0cm, fill=white},
  bulI/.style ={bul, draw=cIn,   fill=cIn!7},
  bulP/.style ={bul, draw=cProc, fill=cProc!8},
  bulL/.style ={bul, draw=cLLM,  fill=cLLM!8},
  bulG/.style ={bul, draw=cGT,   fill=cGT!9},
  bulR/.style ={bul, draw=cErr,  fill=cErr!7},
  bulGy/.style={bul, draw=cGy,   fill=black!4},
  bloc/.style={draw=cGy!60, dashed, rounded corners=3pt, inner sep=6pt, fill=cBg},
  blocT/.style={font=\scriptsize\bfseries, text=cIn},
  FL/.style={-{Latex[length=3mm]}, line width=1pt, draw=cIn!70},
  fl/.style={-{Latex[length=2.2mm]}, thick},
  fld/.style={-{Latex[length=2.0mm]}, thick, dashed},
  et/.style={font=\tiny\itshape, text=cComment, inner sep=1.5pt},
}

% ---------- Style code ----------
\lstset{
  basicstyle=\ttfamily\scriptsize, backgroundcolor=\color{cCode},
  frame=single, framerule=0pt, rulecolor=\color{cCode},
  xleftmargin=6pt, xrightmargin=3pt, framexleftmargin=5pt,
  breaklines=true, columns=fullflexible, keepspaces=true,
  showstringspaces=false, tabsize=2, aboveskip=3pt, belowskip=3pt,
  commentstyle=\color{cComment}\itshape, stringstyle=\color{cStr},
  literate=%
    {é}{{\'e}}1 {è}{{\`e}}1 {ê}{{\^e}}1 {à}{{\`a}}1 {â}{{\^a}}1
    {î}{{\^\i}}1 {ï}{{\"\i}}1 {ô}{{\^o}}1 {û}{{\^u}}1 {ù}{{\`u}}1
    {ç}{{\c c}}1 {É}{{\'E}}1 {«}{{\guillemotleft}}1 {»}{{\guillemotright}}1
    {’}{{'}}1 {–}{{-}}1 {…}{{...}}1,
}
\lstdefinestyle{bash}{language=bash}
\lstdefinestyle{json}{language=,moredelim=[s][\color{cIn}\bfseries]{"}{":}}
\captionsetup{font=small,labelfont=bf,skip=3pt}
\setlength{\parskip}{2pt}

\begin{document}
\pagestyle{empty}

\section*{Annexe A — Évaluation multi-LLM du mapping (workflow n8n \emph{sans} MCP) et métriques}

Cette annexe décrit comment les LLM sont évalués sur la tâche de \textbf{mapping} :
à partir d'un échantillon de télémétrie (produit selon l'annexe précédente),
chaque modèle doit rendre le \textbf{triplet ATT\&CK} attendu — tactique,
technique, sous-technique — dans la bonne matrice. Deux briques : un
\textbf{workflow n8n} qui interroge les modèles et journalise leurs réponses, puis
le script \lstinline|calculer_metriques.py| qui confronte ces réponses à la
vérité-terrain. Le corpus compte \textbf{42 échantillons}
(\texttt{S000.jsonl}\dots\texttt{S0041.jsonl}).

% =====================================================================
\subsection*{A.1\quad Le workflow n8n, présenté par blocs}
% =====================================================================
Le workflow se lit \textbf{de gauche à droite} (Figure~\ref{fig:wf}). Chaque
\textbf{bloc} regroupe plusieurs nœuds n8n (les bulles) ; on décrit ici le rôle
des blocs, pas de chaque bulle.

\begin{figure}[ht]\centering
\resizebox{\linewidth}{!}{%
\begin{tikzpicture}[x=1cm,y=1cm]
% ---- bulles (definies d'abord) ----
\node[bulI] (n1a) at (0,0.5)  {Déclencheur};
\node[bulI] (n1b) at (0,-0.4) {Configuration};
\node[bulI] (n2a) at (3.0,0.5)  {Lire le corpus};
\node[bulI] (n2b) at (3.0,-0.4) {Extraire};
\node[bulP] (n3)  at (6.0,0.05) {Préparer les itérations};
\node[bulGy](n4a) at (9.6,1.85) {Auth Vertex (jeton Google)};
\node[bulL] (n4b) at (9.6,0.95) {Préparer + Appel du modèle};
\node[bulP] (n4c) at (9.6,0.05) {Jugement (statut)};
\node[bulR] (n4d) at (9.6,-0.85){Reprise (backoff)};
\node[bulG] (n4e) at (9.6,-1.75){Écriture (ajout) + pause};
\node[bulG] (n5)  at (12.9,0.05) {Synthèse de campagne};
\node[bulG] (res) at (9.6,-3.3) {resultats\_\_<modele>.jsonl};
% ---- conteneurs de blocs, EN ARRIERE-PLAN (pour ne pas masquer les bulles) ----
\begin{scope}[on background layer]
\node[bloc,fit=(n1a)(n1b),label={[blocT]above:1 · Réglages}] (B1) {};
\node[bloc,fit=(n2a)(n2b),label={[blocT]above:2 · Corpus (42)}] (B2) {};
\node[bloc,fit=(n3),label={[blocT]above:3 · Dépliage}] (B3) {};
\node[bloc,fit=(n4a)(n4b)(n4c)(n4d)(n4e),
      label={[blocT]above:4 · Boucle — un appel à la fois}] (B4) {};
\node[bloc,fit=(n5),label={[blocT]above:5 · Fin}] (B5) {};
\end{scope}
% ---- fleches internes au bloc 4 ----
\draw[fl] (n4a)--(n4b); \draw[fl] (n4b)--(n4c);
\draw[fl] (n4c)--node[et,right]{ok} (n4e);
\draw[fl] (n4c.west) to[out=180,in=180] node[et,left]{erreur} (n4d.west);
\draw[fl] (n4d.east) to[out=0,in=0] node[et,right]{réessaie} (n4b.east);
% ---- fleches entre blocs ----
\draw[FL] (B1.east)--(B2.west);
\draw[FL] (B2.east)--(B3.west);
\draw[FL] (B3.east)--(B4.west);
\draw[FL] (B4.east)--node[et,above]{fin de boucle} (B5.west);
% ---- boucle sur les iterations ----
\draw[fl,draw=cGy] (B4.north east) to[out=55,in=125]
      node[et,above]{itération suivante} (B4.north west);
% ---- sortie fichier ----
\draw[fl,draw=cGT] (B4.south)--(res.north);
\draw[fld] (res.east) -- ++(3.4,0)
      node[right,font=\tiny\itshape,text=cComment,align=left]
      {vers \lstinline|calculer_metriques.py|\\ (A.3)};
\end{tikzpicture}}
\caption{Workflow n8n d'évaluation, de gauche à droite. Le bloc~4 (la boucle)
traite un appel à la fois : authentification, appel du modèle, jugement, reprise
sur erreur, écriture puis pause, avant l'itération suivante. La vérité-terrain
n'entre jamais dans ce workflow.}
\label{fig:wf}
\end{figure}

\noindent\textbf{1 · Réglages} — le point d'entrée et le \emph{seul} nœud à
configurer (clés, modèle, paramètres). \textbf{2 · Corpus} — lecture des 42
échantillons. \textbf{3 · Dépliage} — construction de la liste des appels
(\emph{échantillon $\times$ run $\times$ modèle}). \textbf{4 · Boucle} — le cœur :
un appel à la fois, avec authentification (Vertex pour les modèles Google), appel
du modèle, jugement de la réponse, reprise en cas d'erreur transitoire, puis
écriture au fil de l'eau et pause avant l'itération suivante. \textbf{5 · Fin} —
synthèse de campagne. Le workflow appelle les modèles \textbf{en direct par API
HTTP, \emph{sans} MCP} : aucun outil n'est fourni, le modèle répond à partir du
seul échantillon, ce qui mesure sa connaissance propre du référentiel.

\paragraph{Paramètres choisis et justification.}
\begin{center}\small
\setlength{\tabcolsep}{4pt}\renewcommand{\arraystretch}{1.12}\begin{tabular}{@{}p{4.2cm}p{2.0cm}p{9.3cm}@{}}
\toprule
\textbf{Paramètre} & \textbf{Valeur} & \textbf{Pourquoi} \\
\midrule
\lstinline|runs_par_question| & 3 & chaque échantillon est rejoué plusieurs fois pour \emph{mesurer la stabilité} des réponses (un modèle peut varier même à température nulle). \\
\lstinline|temperature| & 0 & réponses les plus déterministes et reproductibles possible. \\
\lstinline|limite_samples| & 0 (= 42) & traite tout le corpus ; on met 4–5 uniquement pour un test de coût. \\
\lstinline|max_tokens_sortie| & 32768 & plafond \emph{identique} pour tous, assez haut pour n'être jamais atteint (\texttt{meta.fin\_raison} prouve après coup qu'aucune réponse n'est tronquée), tout en bornant le coût d'une génération qui s'emballe. Obligatoire chez Anthropic. \\
\lstinline|MODELE_FORCE| (1 modèle / exécution) & \texttt{""} ou un modèle & permet de lancer plusieurs campagnes \emph{en parallèle}, chacune écrivant son propre fichier — donc aucune collision. \\
\lstinline|pause_entre_appels_s| & 8 (défaut) & espacement anti-quota entre deux appels réussis ; surchargé par modèle (p.\,ex.\ Moonshot 21\,s, limité à 3 req/min). \\
Backoff (reprise) & $\le$ 60\,s, \lstinline|max_tentatives| & sur 429 / 5xx / réseau : on respecte d'abord le délai demandé par le serveur (\texttt{Retry-After}\dots), sinon backoff exponentiel avec \emph{jitter}. \\
Appel en \texttt{fullResponse} + \texttt{neverError} & — & on récupère le code HTTP et le corps \emph{même en erreur} : c'est ce qui rend la reprise possible et permet de séparer erreurs du modèle et de l'infrastructure. \\
Sortie : 1 fichier / modèle, écriture en ajout & — & robustesse (un plantage ne perd que l'appel courant) et parallélisme (pas de collision entre campagnes). \\
\lstinline|exclure_des_metriques| & drapeau & marque les pannes d'infrastructure (quota, annulation, timeout) pour qu'elles ne comptent pas comme des erreurs du modèle. \\
\bottomrule
\end{tabular}
\end{center}

% =====================================================================
\subsection*{A.2\quad Ce que produit le workflow}
% =====================================================================
Chaque appel donne \textbf{une ligne JSON} ajoutée au fichier du modèle
(\texttt{resultats/resultats\_\_<modele>.jsonl}) : la réponse du modèle (le triplet
ATT\&CK), son \textbf{statut} (\texttt{ok}, \texttt{format\_non\_respecte},
\texttt{erreur\_quota\_429}\dots), et des méta-données d'exécution (latence, jetons,
numéro de run). Ces fichiers sont l'entrée du calcul des métriques.

% =====================================================================
\subsection*{A.3\quad Calcul des métriques : \texttt{calculer\_metriques.py}}
% =====================================================================
Le script confronte les fichiers de résultats à la vérité-terrain et produit le
rapport d'évaluation. Il télécharge puis met en \textbf{cache} le référentiel
ATT\&CK (Enterprise \emph{et} ICS) pour \emph{juger} chaque identifiant prédit.

\begin{figure}[ht]\centering
\begin{tikzpicture}[x=1cm,y=1cm,node distance=3mm]
\node[bulG,text width=3.2cm]  (r)  at (0,0.9)  {\texttt{resultats\_\_*.jsonl}\\ \scriptsize (workflow)};
\node[bulI,text width=3.2cm]  (v)  at (0,0)    {\texttt{verite\_terrain.jsonl}\\ \scriptsize (vérité-terrain)};
\node[bulGy,text width=3.2cm] (m)  at (0,-0.9) {Référentiel MITRE\\ \scriptsize (cache, 2 matrices)};
\node[bulP,text width=2.9cm]  (f)  at (4.3,0)  {Filtrage + dédup\\ \scriptsize écarte erreurs d'infra};
\node[bulP,text width=3.0cm]  (e)  at (8.2,0)  {Étiquetage 3 niveaux\\ \scriptsize VP/FP/FN/HALLU/OBSOLÈTE};
\node[bulP,text width=3.0cm]  (met)at (8.2,-1.9){Métriques par niveau\\ \scriptsize P, R, F1 (strict \& tolérant)};
\node[bulG,text width=3.0cm]  (rap)at (4.3,-1.9){\textbf{Rapport}\\ \scriptsize tableaux + classements};
\draw[fl] (r.east)  to[out=0,in=180] (f.west);
\draw[fl] (v.east)  -- (f.west);
\draw[fl] (m.east)  to[out=0,in=180] (f.west);
\draw[fl] (f) -- (e);
\draw[fl] (e) -- (met);
\draw[fl] (met) -- (rap);
\end{tikzpicture}
\caption{Chaîne du calcul des métriques : entrées $\to$ filtrage $\to$ étiquetage
aux trois niveaux $\to$ métriques et classements.}
\label{fig:metr}
\end{figure}

\textbf{Filtrage préalable.} Les itérations abandonnées pour cause
d'infrastructure (quota, annulation, timeout, réseau — drapeau
\lstinline|exclure_des_metriques|) sont \textbf{écartées} : les compter en faux
négatifs attribuerait au modèle une panne de la plomberie. Comme les fichiers sont
écrits en ajout, les doublons éventuels sont dédupliqués (on garde la dernière
occurrence de chaque couple \emph{modèle $\times$ échantillon $\times$ exécution}).

\textbf{Étiquetage}, appliqué à chaque exécution et à \textbf{chaque niveau}
(tactique, technique, sous-technique) :
\begin{center}\small
\begin{tabular}{@{}ll@{}}
\toprule
\textbf{VP} & l'identifiant attendu est prédit \\
\textbf{FP} & un identifiant existant dans la version figée, mais incorrect \\
\textbf{FN} & le bon identifiant n'a pas été identifié (prédiction fausse \emph{ou} abstention) \\
\textbf{HALLUCINATION} & identifiant absent de \emph{toute} version du référentiel (compté en FP) \\
\textbf{OBSOLÈTE} & identifiant révoqué dont le remplaçant est l'attendu \\
\bottomrule
\end{tabular}
\end{center}

\noindent Une réponse rejetée (format non respecté, erreur d'API) est traitée
comme une \textbf{abstention} aux trois niveaux. Cas de la \textbf{sous-technique}:
si la technique attendue n'en a pas, « aucune » est un VP et affirmer une
sous-technique est une \emph{sur-spécification} (FP) ; si elle en a une,
l'abstention est un FN. Le F1 de sous-technique n'est calculé que sur les
échantillons dont la technique \emph{possède} une sous-technique (22 sur 42), pour
éviter qu'un modèle muet obtienne un F1 flatteur grâce aux échantillons terminaux.

\textbf{Métriques}, calculées \textbf{séparément à chaque niveau} (jamais de F1
global) : $\text{Précision}=\tfrac{VP}{VP+FP}$, $\text{Rappel}=\tfrac{VP}{VP+FN}$,
$F1=\tfrac{2PR}{P+R}$. Un modèle qui s'abstient perd du rappel sans perdre de
précision ; un modèle qui répond faux perd les deux. Le \textbf{taux d'abstention}
complète la lecture, et les F1 sont \textbf{recalculés hors rejets} pour séparer la
compétence de mapping de la conformité au format. Chaque F1 est calculé
\textbf{deux fois} : \emph{strict} (les obsolètes comptent comme erreurs, métrique
principale) et \emph{tolérant} (les obsolètes comptent comme VP) ; l'écart mesure
la part d'erreurs due à l'obsolescence des connaissances du modèle.

\textbf{Diagnostics et classements.} Le rapport ajoute des attributs de diagnostic
(mauvaise matrice, taux d'hallucination, d'obsolescence, d'abstention,
sur/sous-spécification) et une vue \emph{fiabilité/coût} (stabilité — part des
échantillons dont toutes les exécutions rendent le même triplet —, latence p50/p95,
jetons moyens, efficacité). Deux classements en découlent : \textbf{qualité du
mapping} (mauvaise matrice, puis hallucination, puis F1 technique~$>$~tactique~$>$~sous-technique,
puis obsolescence et abstention) et \textbf{fiabilité} (stabilité, puis latence,
puis coût). Le script prévoit aussi une phase « avec serveur MCP » ; la présente
campagne étant \emph{sans} MCP, ces indicateurs restent vides.

\begin{lstlisting}[style=bash]
# usage : un fichier de resultats par modele + la verite de terrain
uv run calculer_metriques.py "resultats/resultats__*.jsonl" --verite verite_terrain.jsonl --sortie rapport
\end{lstlisting}

\end{document}
